% 6.1 章 · 二维综合仿真
% 这一章的笔记把前五部分的零件装成三个完整求解器，并把"能出图"与"算对了"之间的
% 证据链拆成四页：方腔的残差 vs 离散误差、Ghia 基准的抄写地板、Gray--Scott 的
% 图灵判据（含三种"看起来像图灵"的对照）、干底溃坝的误差窗口与质量账本。
% 所有数字都来自 tests/6.1/ 的断言记录（asserts_table.tex 由断言报告生成）。

\numlabchapter{二维综合仿真}{6.1}

\section{这一章要解决什么}

前面各部分已经有了离散算子、时间推进、线性求解器、谱方法与优化器。本章不再引入
新的数学工具，而是把它们组装成三个**完整的应用求解器**，然后只问一个问题：
\[
  \text{“图看起来对”} \quad\Longrightarrow\quad \text{“这个结论有证据吗？”}
\]
四个可视化分别对应四条独立证据链：

\begin{enumerate}
  \item \textbf{方腔}（\texttt{cavity\_vorticity}）：顶盖驱动、$\psi$--$\omega$ 形式、
        DST 直接解 Poisson、红黑 SOR 扫涡量。残差 $10^{-7}$ 与 Ghia 中心线
        RMSE $0.2271$ 同时成立——前者量离散方程，后者量离散误差。
  \item \textbf{Ghia 基准}（\texttt{ghia\_benchmark}）：把数值解插到
        Ghia--Ghia--Shin (1982) 的原表节点上，分别对拍 $u(y)$ 与 $v(x)$；
        Re=100 的四网格 RMSE 为 $0.02024,0.01211,0.00870,0.00603$，实测收敛阶
        $1.10$。基准本身有 $10^{-2}$ 量级的转载分歧，页面把这条地板画出来。
  \item \textbf{Gray--Scott}（\texttt{turing\_pattern}）：真图灵预设的
        $\varphi^\ast=53.8988$、$\lambda^\ast=0.015017$；det 的顶点
        $\varphi=62.7495$ 只给 $0.013392$，低 $10.8\%$。经典 spots 的
        $\lambda^\ast=-0.035$ 仍长出振幅 $0.729$ 的斑图；
        stripes 的底态存在但 $\lambda(0)=+0.00539$（均匀失稳）；
        uniform 的 $\lambda(\varphi)=-D\varphi^2$ 只有纯扩散衰减。
  \item \textbf{干底溃坝}（\texttt{dam\_break}）：Rusanov 通量 + 显式推进，
        与 Ritter 精确解对拍。$t=0.06$ 时扇内 L1 $=2.02\times10^{-2}$，
        全域 L1 $=2.21\times10^{-2}$，比值 $1.093$；$n=100\to400$ 的扇内 L1
        从 $3.04\times10^{-2}$ 降到 $1.25\times10^{-2}$。质量漂移 $9.9\times10^{-9}$
        来自正性地板注水，不是守恒格式的反证。
\end{enumerate}

记号：$h=1/(n-1)$ 是方腔网格间距；$\omega=\partial_xv-\partial_yu$ 是涡量，
$\psi$ 是流函数；$U=1$ 是顶盖速度；$\mathrm{RMSE}$ 是数值解与基准在基准节点上的
均方根差；$\varphi$ 是连续波数（弧度/单位长），$z=\varphi^2$；
$D_u,D_v,F,k$ 是 Gray--Scott 的扩散系数与喂料/杀灭率；
$c_L=\sqrt{g h_L}$ 是干底溃坝的左行波速，$\xi=(x-x_0)/t$。

\section{方腔：残差不是误差}\label{sec:cavity}

\subsection{流函数--涡量形式与两个边界条件}

不可压方腔里，速度由流函数给出：
\[
  u=\frac{\partial\psi}{\partial y},\qquad
  v=-\frac{\partial\psi}{\partial x},\qquad
  \Delta\psi=-\omega .
\]
第二个式子是 Poisson 方程，配齐四壁的 $\psi=0$；涡量方程是
\[
  \omega_t+u\omega_x+v\omega_y=\frac1{Re}\Delta\omega .
\]
壁面涡量不能直接由 $\psi$ 的内部差分得到（会损失一阶精度、并且符号容易写反），
实现用 Thom 公式
\[
  \omega_w=-\frac{2\psi_{\text{邻}}}{h^2}-\frac{2U_\tau}{h},
\]
其中底/左/右壁 $U_\tau=0$，顶盖 $U_\tau=1$。顶盖沿 $+x$ 运动时顶壁涡量约为
$-2/h$；符号写反会得到一个**镜像**但仍然光滑的解——$n=129$、$Re=100$ 时主涡心会
跑到 $x=0.383$ 而不是 Ghia 的 $0.6172$，而残差曲线依然漂亮。这是本章最值得
先讲的一条“可测量陷阱”。

\subsection{DST：Poisson 与 Helmholtz 都被对角化}

\begin{theorem}[五点 Laplace 的 DST 对角化]\label{thm:dst}
在 $m=n-2$ 个内部点上，记
$\lambda_{pq}=-\frac{4}{h^2}\bigl(\sin^2\frac{p\pi}{2(m+1)}
+\sin^2\frac{q\pi}{2(m+1)}\bigr)$。则五点 Laplace 的特征向量为
$(\sin\frac{ip\pi}{m+1}\sin\frac{jq\pi}{m+1})_{ij}$，特征值为
$\lambda_{pq}$；因此
\[
  \Delta_h\psi=-\omega,\qquad (I-a\Delta_h)\psi=f
\]
都可以用二维 DST 在 $O(n^2\log n)$ 内直接求解，而不需要迭代。
\end{theorem}

\begin{proof}
把五点模板写成 $L_h=L_x\otimes I+I\otimes L_y$，其中 $L_x$ 是一维 Dirichlet
二阶差分。一维 $L_x$ 的特征向量是 $\sin(ip\pi/(m+1))$，特征值为
$-4\sin^2(p\pi/(2(m+1)))/h^2$；张量积的特征值相加，特征向量取乘积。于是
Poisson 在变换域里逐点除以 $\lambda_{pq}$，Helmholtz 逐点除以
$1-a\lambda_{pq}$（$a=0$ 退化为 Poisson）。实现中对四个角点的边界行/列单独处理，
内部点数与变换长度一致，所以没有额外填充。
\end{proof}

\begin{proof}[制造解实测]
取 $\psi_{\text{exact}}=\sin(\pi x)\sin(\pi y)$，五点 Laplace 得到 $f$ 后再用
同样的 DST 求解。实测 $\max|\Delta_h\psi+\omega|=4.885\times10^{-14}$，
Helmholtz 的制造解误差 $\max|\psi_H-\psi_{\text{exact}}|=1.110\times10^{-15}$。
这两条是 A 类断言的“解析靶子”：它们证明离散算子与变换实现没有互相抵消的错。
\end{proof}

由 $\psi$ 求速度时，$u=\psi_y$、$v=-\psi_x$ 用同源的中心差分；离散散度
$\max|u_x+v_y|=8.882\times10^{-16}$。这条恒等式不是“看起来满足”，而是网格内
每个点都成立的代数事实。

\subsection{耦合迭代与相对残差}

外层用红黑 SOR 扫涡量方程，内层每个 sweep 用一次 DST Poisson；收敛判据是
\[
  R^{(k)}=\frac{\max|R_{\text{disc}}|}
  {\max|u|/h\cdot\max|\omega|}.
\]
实现记录每一次 sweep 的 $R$，最后把 $\log_{10}R$ 交给前端。默认参数
$Re=1000,n=41$ 时，1087 次 SOR + 1087 次 Poisson 把 $R$ 压到
$1.00\times10^{-7}$（tol $=10^{-7}$）。这**只说明**离散稳态方程被解得足够好；
与 Ghia 的中心线 RMSE $0.2271$（$u$ 剖面 $0.1833$、$v$ 剖面 $0.2637$）相比，
中间差了四个数量级。

\begin{theorem}[残差与误差是两件事]\label{thm:residual}
设 $L_h\psi_h=\omega_h$ 是离散方程，$\psi$ 是连续解。相对残差
$R=\|L_h\psi_h-f_h\|/\text{尺度}$ 可以在 $h\to0$ 前任意小（只要解方程够狠），
而离散误差 $\|\psi_h-\psi\|\sim C h^p$ 由格式与网格决定，不会被迭代误差压低到
$Ch^p$ 以下。因此“残差小”既不蕴含“误差小”，也不能替代网格收敛研究。
\end{theorem}

\begin{proof}
把误差分解为 $\psi_h-\psi=(\psi_h-\psi_h^\ast)+(\psi_h^\ast-\psi)$，其中
$\psi_h^\ast$ 是离散方程的精确解。第一项由迭代误差控制，第二项由截断误差控制。
残差只直接约束第一项；第二项需要与基准/制造解比较，或做网格加密。
$Re=1000,n=33$ 的实测是干净的负对照：残差 $9.89\times10^{-8}$，但
Ghia RMSE $0.2413$；同一套代码在 $Re=100,n=49$ 上给出 $9.91\times10^{-8}$ 与
$0.0060$。残差相同，误差相差 40 倍，说明差异只能来自离散误差。
\end{proof}

\subsection{时间步的两个来源}

如果要把稳态求解器改成显式推进，两个稳定性限制会同时出现：
\begin{align}
  \Delta t_{\text{diff}}&=\frac{h^2Re}{4},\\
  \Delta t_{\text{CFL}}&=\frac{c_{\text{cfl}}h}{U_{\max}}.
\end{align}
交叉点满足 $\Delta t_{\text{diff}}=\Delta t_{\text{CFL}}$，即
\[
  Re_\times=\frac{4c_{\text{cfl}}(n-1)}{U_{\max}}.
\]
$c_{\text{cfl}}=0.8$、$U_{\max}=1$ 时 $Re_\times=3.2(n-1)$：$n=33$ 给 $102.4$，
$n=41$ 给 $128$，$n=49$ 给 $153.6$。在 $Re=1000$ 的默认页上，$n=41$ 的
$\Delta t_{\text{CFL}}=0.0200$ 小于 $\Delta t_{\text{diff}}=0.1563$，
由 CFL 绑定；$Re=50$ 时反过来，由扩散绑定。
“Re 越大越难算”这句话在交叉点两侧方向相反。

另一条容易忽略的限制是中心差分的格点雷诺数：
\[
  Re_h=Re\,|u|\,h=\frac{Re}{n-1}.
\]
五点中心对流的对角占优条件是 $Re_h<2$，即 $Re<2(n-1)$。$n=41$ 时阈值
$Re=80$，而 $Re=1000$ 的 $Re_h=25$——中心格式会先失稳而不是先变准；
实现因此默认一阶迎风。迎风带来 $O(h)$ 的数值扩散，也是 Ghia 曲线在
$Re=1000$ 粗网格上误差仍有 $0.2$ 的来源。

\begin{center}
\begin{tabular}{lrrrr}
\toprule
量 & $Re=100,n=49$ & $Re=1000,n=41$ & 解释 \\
\midrule
SOR 残差 & $9.9\times10^{-8}$ & $1.0\times10^{-7}$ & 离散方程离差 \\
Ghia RMSE & $0.0060$ & $0.2271$ & 与基准的差 \\
$Re_h$ & $2.04$ & $25.00$ & 中心格式阈值 $2$ \\
时间步绑定 & 扩散 & CFL & $Re_\times=3.2(n-1)$ \\
\bottomrule
\end{tabular}
\end{center}

\section{Ghia 基准：把“基准”本身也当成误差源}\label{sec:ghia}

\subsection{表、方向与端点}

Ghia 表给的是两条中心线上的离散节点：竖直中心线 $u(y)$ 与水平中心线 $v(x)$，
每条 17 个点。实现把**数值解插值到基准节点**，而不是把基准插到数值网格上：

\begin{itemize}
  \item 后者会把基准的局部极值抹平，误差曲线看起来更漂亮；
  \item 数值解在基准节点上的插值误差是 $O(h^2)$，远小于一阶迎风的 $O(h)$；
  \item 端点必须精确：$u(0)=u(1)=0$（除顶盖 $u(1)=1$）、$v(0)=v(1)=0$，
        用端点常数外插而不是线性外推。
\end{itemize}

实测三列（$Re=100,400,1000$）的端点严格相等、采样严格递增；主涡中心随 Re
单调移动：$(0.6172,0.7344)\to(0.5547,0.6055)\to(0.5313,0.5625)$。
$Re=1000$ 还存在左下/右下两个二次涡。

\subsection{抄写地板}

不同转载的 Ghia 表在 $Re=1000$ 的 $v$ 剖面上，$x=0.9063$ 处就有
$-0.51550$ 与 $-0.51500$ 两种写法。把容差压到 $10^{-3}$ 的“完美对拍”量的是
转载版本，不是数值误差。页面固定报 \texttt{transcription\_floor=1e-2}，并把这条虚线
画在误差图上；当 RMSE 与地板同量级时，读者应该读“已经到抄写精度”，而不是
“格式只有这么准”。

\subsection{网格收敛阶}

用 $h=1/(n-1)$、四档 $n=17,25,33,49$ 做 log--log 拟合，$Re=100$ 的总 RMSE
为 $0.02024,0.01211,0.00870,0.00603$，斜率 $1.10$；$u$ 剖面 $0.64$、$v$ 剖面
$2.10$。一阶迎风的理论阶是 1，实测在 $(0.5,2.5)$ 的宽窗内，但**不能把 $v$ 的
2.10 读成二阶精度**：$v$ 的误差更小、更接近抄写地板，四点的局部斜率容易被
基准本身的截断带偏。$Re=1000$ 时总斜率 $0.088$、$u$ 0.19、$v$ 0.19，
粗网格上数值扩散主导，曲线接近平线；这不是代码错，而是离散误差大。

\begin{center}
\begin{tabular}{rrrrr}
\toprule
$n$ & 总 RMSE & $u$ 剖面 & $v$ 剖面 & 备注 \\
\midrule
17 & 0.02024 & 0.01677 & 0.02320 & 粗网格 \\
25 & 0.01211 & — & — & 下降 \\
33 & 0.00870 & — & — & 下降 \\
49 & 0.00603 & 0.00821 & 0.00231 & 最优点 \\
\bottomrule
\end{tabular}
\end{center}

\section{Gray--Scott：图灵判据不是“图案好看”}\label{sec:turing}

\subsection{方程与非平凡稳态}

单位正方形周期域上的 Gray--Scott 系统
\begin{align}
  u_t&=D_u\Delta u-uv^2+F(1-u),\\
  v_t&=D_v\Delta v+uv^2-(F+k)v .
\end{align}
均匀稳态满足 $v=(F+k)/u$。代入 $u$ 方程，得到
\[
  u_\pm=\frac{1\pm\sqrt{1-4(F+k)^2/F}}{2},
\]
并要求判别式非负。两支中只有 $u_-$（小 $u$、大 $v$）给出
$\det J=(F+k)(v^2-F)>0$；$u_+$ 是 $\det J<0$ 的鞍点。
turing 预设 $(F,k)=(0.03,0.045)$ 的底态是 $(0.25,0.30)$，$\det J=0.0045$；
spots 预设 $(0.035,0.060)$ 的判别式为负，没有非平凡稳态，实现退回平凡态
$(1,0)$。

\subsection{线性化与图灵条件}

在底态 $J=\begin{pmatrix}f_u&f_v\\g_u&g_v\end{pmatrix}$ 上加模态
$\exp(\lambda t+i\varphi x)$，得到
\[
  M(\varphi)=J-\varphi^2\operatorname{diag}(D_u,D_v),\qquad
  \lambda_{\max}(\varphi)=\frac{\operatorname{tr}M}
  {2}+\frac{\sqrt{\operatorname{tr}^2M-4\det M}}{2}.
\]
图灵失稳要求：
\begin{enumerate}
  \item 底态无扩散稳定：$\operatorname{tr}J<0$、$\det J>0$；
  \item 存在 $\varphi>0$ 使 $\lambda_{\max}(\varphi)>0$；
  \item 边界 $\varphi=0$ 仍然稳定，即 $\lambda(0)<0$。
\end{enumerate}

\begin{theorem}[图灵带与最不稳定波数]\label{thm:turing}
记 $z=\varphi^2$，
\[
  \det M(z)=D_uD_vz^2-(D_ug_v+D_vf_u)z+\det J .
\]
在 $\operatorname{tr}J<0,\det J>0$ 的前提下，图灵失稳等价于
$D_ug_v+D_vf_u>2\sqrt{D_uD_v\det J}$，且失稳区间是某个 $z>0$ 的开区间。
最不稳定波数由
$\frac{d}{dz}\bigl(\operatorname{tr}M+\sqrt{\operatorname{tr}^2M-4\det M}\bigr)=0$
给出，一般**不是** $\det M$ 的顶点
$z_v=(D_ug_v+D_vf_u)/(2D_uD_v)$。
\end{theorem}

\begin{proof}
$\lambda_{\max}(z)$ 是 $M$ 的最大特征值。$\det M(z)$ 在 $z_v$ 取最小，但
$\operatorname{tr}M(z)=f_u+g_v-(D_u+D_v)z$ 随 $z$ 线性下降；当特征值为实数时
$\lambda_{\max}=(\operatorname{tr}+\sqrt{\operatorname{tr}^2-4\det})/2$，
两项的竞争把最大点推到 $z_v$ 左侧。对 $z$ 求导并解二次方程即得驻点；
端点 $z=0$ 与复特征值分支单独比较，取全局最大。turing 预设的实测是
$\varphi^\ast=53.8988$、$\lambda^\ast=0.015017$，而 det 顶点 $\varphi=62.7495$
给 $0.013392$，低 $10.8\%$；只信 det 顶点会系统性偏大波长。
\end{proof}

\paragraph{四种预设的正确读法}

\begin{center}
\begin{tabular}{llll}
\toprule
预设 & 底态 & $\lambda(0)$ & 结论 \\
\midrule
turing & $(0.25,0.30)$，$\det>0$ & $-0.0225$ & 真图灵带 $\varphi^\ast=53.90$ \\
stripes & $(0.404,0.210)$ & $+0.00539$ & 均匀模失稳，不是图灵带 \\
spots & 无（退回平凡态） & $-0.035$ & 自复制瞬态 \\
waves & 无（退回平凡态） & $-0.014$ & 自复制瞬态 \\
uniform & 平凡态 $(1,0)$ & $0$ & 纯扩散 $\lambda=-D\varphi^2$ \\
\bottomrule
\end{tabular}
\end{center}

spots 的负对照最容易骗人：$\lambda^\ast=-0.035$ 却仍能在 4000 步后长出
振幅 $0.729$、16 个斑块的图案。图案来自激励介质的自复制瞬态，不是线性图灵
失稳；把“有斑图”当“有正特征值”会把整类 Gray--Scott 图案误判。
stripes 的底态存在，但 $g_v=F+k>0$ 且 $\operatorname{tr}J>0$：它在 $\varphi=0$
就已经均匀失稳，也不是图灵带。uniform 的反应项为零，色散关系
$\lambda(\varphi)=-D\varphi^2$ 全是衰减，没有增长模。

\subsection{显式时间步与模板误差}

FTCS 的五点模板在 Fourier 域里的符号是 $[-8/\Delta x^2,0]$，放大因子
$1-8D\Delta t/\Delta x^2$ 必须落在 $[-1,1]$，因此
\[
  \Delta t\le\frac{\Delta x^2}{4\max(D_u,D_v)}.
\]
实现默认再乘安全系数 $0.5$。$n=128$、turing 预设的 $\Delta x=1/128$，
上限 $0.190735$，安全值 $0.0953674$；请求的 $0.25$ 被夹到安全值，帧里的
\texttt{dt\_clamped=1}、\texttt{dt\_requested=0.25}、\texttt{dt\_limit=0.0953674} 三者同时存在，页面
不把“请求值”冒充“实际值”。\texttt{clamp=False} 时直接抛 \texttt{ValueError}。

空间模板的有效波数是
\[
  \varphi_{\text{eff}}=\frac{2}{\Delta x}\sin\left(\frac{\varphi\Delta x}{2}\right),
\]
把 $\varphi_{\text{eff}}$ 代入解析增长率就得到离散增长率。误差随 $\Delta x^2$
下降：$n=64,128,256,512$ 的相对误差
$9.32\times10^{-3},3.24\times10^{-3},8.68\times10^{-4},2.21\times10^{-4}$，
实测阶 $1.81$。这是**空间模板误差**，不是完整时间推进误差；页面把两件事分开。

\subsection{质量与波长}

扩散项在周期域上逐点和为零，因此
\[
  \frac{d}{dt}\sum_{ij}u_{ij}=\Delta t\sum_{ij}
  \bigl(-u_{ij}v_{ij}^2+F(1-u_{ij})\bigr).
\]
五步实测的 $\Delta\Sigma u$ 与右端预测相对差为 $0$。Gray--Scott 不守恒；
质量变化本身不是错误。

图灵饱和图案的主波长由 $\varphi^\ast$ 选出：理论 $2\pi/\varphi^\ast=0.1166$，
$n=128$、10000 步实测 $0.1253$（差 $7.5\%$），振幅 $0.302$、96 个局部极大。
有限盒、初始随机种子与非线性饱和都让实测偏离解析值；页面把“主波长无定义”
和“波长恰好为 0”严格分开：uniform 的 \texttt{wavelength\_valid=0}，帧里存 0 只是占位，
侧栏写“无空间结构”，不写“波长 0”。

\section{干底溃坝：窗口、守恒与 CFL}\label{sec:dam}

\subsection{浅水方程与 Rusanov 通量}

一维浅水方程
\[
  \partial_t\begin{pmatrix}h\\hu\end{pmatrix}
  +\partial_x\begin{pmatrix}hu\\hu^2+\tfrac12gh^2\end{pmatrix}=0 .
\]
有限体积的 Rusanov 通量
\[
  \hat F_{i+1/2}=\frac12(F_L+F_R)-\frac12 a_{i+1/2}(U_R-U_L),
  \qquad a_{i+1/2}=\max(|u_L|+\sqrt{gh_L},|u_R|+\sqrt{gh_R}).
\]
实测相同态通量逐位相等（$\Delta F_h=\Delta F_{hu}=0$）；波速与
$|u|+\sqrt{gh}$ 逐点一致。干底用正性地板 $h=\max(h,10^{-8})$；
这不是“物理正确”，而是让波速有定义、让推进不出现负水深。质量因此会有
$10^{-9}$ 量级的漂移，页面把它写成“地板注入”，不写成“守恒误差”。

\subsection{Ritter 解与窗口}

初始 $h_L=1$、右侧干底，精确解是自相似扇：
\[
  h(\xi)=\frac{(2c_L-\xi)^2}{9g},\qquad
  u(\xi)=\frac23(\xi+c_L),\qquad -c_L<\xi<2c_L,
\]
$\xi=(x-x_0)/t$，$c_L=\sqrt{gh_L}$。扇外 $h=1$（左侧）或 $h=0$（右侧干区），
干锋在 $\xi=2c_L$。$t=0.06$ 时窗口 $x\in[0.312,0.876]$；扇内是光滑区，
全域还包含干锋折点，后者的位置误差不是格式阶的干净估计。

\begin{center}
\begin{tabular}{rrrr}
\toprule
$n$ & 扇内 L1 & 全域 L1 & 全域/扇区 \\
\midrule
100 & $1.91\times10^{-2}$ & $3.13\times10^{-2}$ & 1.75 \\
200 & $1.20\times10^{-2}$ & $2.07\times10^{-2}$ & 1.75 \\
400 & $7.42\times10^{-3}$ & $1.30\times10^{-2}$ & 1.75 \\
\bottomrule
\end{tabular}
\end{center}

扇内 L1 的实测阶为 $0.669$ 与 $0.694$（n=100→200、200→400），在全域上则被
干锋误差污染。帧里 \texttt{fan} 与 \texttt{error\_valid} 两个掩码分别负责“扇区”与“不是地板
占位”；t=0 时 \texttt{n\_steps=0}、\texttt{fan} 全 0、误差全部无效，面板为空而不是画一条假零线。

\subsection{well-balanced 与 CFL}

静水 $h+\text{底高程}=\text{const}$、$u=0$ 时，通量梯度与源项应逐点抵消。
well-balanced 支路实测质量/水面漂移 $1.89\times10^{-16}$；朴素中心源项给
$3.30\times10^{-4}$，相差 $1.7\times10^{12}$ 倍。后者不是“精度差一点”，
而是会在静水上制造假波。

CFL 不是抄来的常数：$cfl=0.8$ 时到达 $t_{\text{end}}$、$h_{\min}=0.426$；
$cfl=2.0$ 时 830 步后 $h_{\max}=\infty$、解非有限。页面把实际
\texttt{max\_cfl} 与步数写进帧，且不把越界的请求悄悄夹回。

\subsection{二维径向溃坝}

二维版本把初始跃变换成圆形水柱，用 Rusanov + 出流边界推进到 $t=0.05$：
$h_{\min}=0.1$、总质量 $0.278418$ 不增、能量 $0.852330$ 不增、maxCFL $0.30$。
二维的“干区不出现负水深”比一维更脆：对角方向的波速、边界出流、正性地板三者
必须同时正确。这条 A 类断言是本章 2D 求解器的入口检查。

\section{四张图分别在说什么}\label{sec:reading}

\paragraph{方腔页}
左上是稳态涡量（红正蓝负的发散色标）与流函数等值线，黑箭头是速度方向
（按全场最大速度归一，只表方向）。左下是 $\log_{10}$ 相对残差；侧栏同时给
SOR 残差与 Ghia RMSE，并说明两者不能互替。红点标主涡，右侧色标由渲染器
自动配；色标与文字之间留了 110\,px 的过道。

\paragraph{Ghia 页}
左上 $u(y)$、右上 $v(x)$，蓝色是数值解、橙点是 Ghia；左下是总/ $u$ / $v$
三条 RMSE 曲线，灰虚线是 $10^{-2}$ 抄写地板。四个网格点是实际求解出来的，
不是拟合出来的；侧栏写实收阶与最优点。

\paragraph{图灵页}
主图是 $v(x,y)$（viridis，固定色标 0 到当前最大值），右上是解析色散
$\lambda(\varphi)$（灰虚是 0，橙虚是 $\varphi^\ast$），左下是五点模板误差
（对数纵轴、灰虚 $O(n^{-2})$）。四种预设的底态、$\lambda(0)$、$\varphi^\ast$
和“是否有正波数增长带”在侧栏逐项列出；uniform 的色标固定 0--1 并明确写
“$v\equiv0$”，不拿 $10^{-9}$ 的假细节充当色标。

\paragraph{溃坝页}
主图是水深剖面：蓝是数值解、橙虚是 Ritter、灰虚是初值、灰带是扇区；
左下是逐点 $|h-h_{\text{exact}}|$（对数），只画 \texttt{error\_valid=1} 的点。
侧栏把扇内 L1、全域 L1、质量漂移、实际 max CFL、末步 dt 分开报。
t=0 的帧没有扇区、没有有效误差，页面写“0 步”和“无有效点”，不画假数据。

\section{本章的“不要被骗”清单}

\begin{itemize}
  \item 残差 $10^{-7}$ 不等于 RMSE 小：$Re=1000,n=41$ 的 RMSE 仍有 $0.2271$。
  \item Ghia 表本身有 $10^{-2}$ 量级的转载分歧；低于它的容差量的是抄写版本。
  \item 数值解要插到基准节点，不能把基准插到数值网格上。
  \item $\lambda^\ast>0$ 不一定是图灵：stripes 的 $\lambda(0)=+0.00539$ 是均匀失稳。
  \item 有斑图不等于图灵：spots 的 $\lambda^\ast=-0.035$ 仍长出振幅 $0.729$ 的斑图。
  \item 没有正波数增长带时，“主波长”无定义；帧里的 0 只是占位，不是测量值。
  \item dt 被夹紧时必须同时显示请求值、实际上限与实际值；不能只报一个好看数。
  \item 五点模板的 $O(\Delta x^2)$ 是空间误差，不是完整时间推进误差。
  \item 干底的质量漂移主要来自正性地板注入，不是守恒性的反证。
  \item 全域 L1 混着干锋位置误差；报格式阶要用扇内窗口。
  \item well-balanced 与朴素源项在静水上差 $10^{12}$ 倍，不是“小误差”。
  \item CFL 越界会真出非有限值；页面不兜底、不静默夹取。
\end{itemize}

% 本文件由 tools/build_book.py 从 dist/reports/assertions.json 生成，请勿手改。
\begin{longtable}{@{}p{0.06\textwidth}p{0.32\textwidth}p{0.22\textwidth}p{0.16\textwidth}@{}}
\caption{章节 6.1 数值断言（生成时间 2026-10-19T16:08:50-0400）}\\
\toprule
类别 & 断言 & 理论值 & 实测值 \\
\midrule
\endfirsthead
\toprule
类别 & 断言 & 理论值 & 实测值 \\
\midrule
\endhead
\bottomrule
\endlastfoot

A & DST Poisson 残差与 Helmholtz 制造解 & max|$\Delta$\_h$\psi$+$\omega$| <=\allowbreak{} 1e-10；max|$\psi$\_H-$\psi$\_exact| <=\allowbreak{} 1e-10 & max|res|=\allowbreak{}4.885e-14，max|$\psi$\_H-$\psi$\_exact|=\allowbreak{}1.110e-15 \checkmark \\
A & 由流函数得到的速度场散度；区间外插值 & 离散散度 <=\allowbreak{} 1e-12(\allowbreak{}五点/\allowbreak{}中心差分同源)；区间外取端点常数 & max|div|=\allowbreak{}8.882e-16；插值 [1.0,\allowbreak{} 1.0,\allowbreak{} 4.0,\allowbreak{} 7.0,\allowbreak{} 7.0] vs [1.0,\allowbreak{} 1.0,\allowbreak{} 4.0,\allowbreak{} 7.0,\allowbreak{} 7.0] \checkmark \\
B & Re=1000/n=33 与 Re=100/n=49 的残差 vs Ghia RMSE & 两组残差都 <1e-6；Re=\allowbreak{}1000 粗网格 RMSE>0.2，而 Re=\allowbreak{}100 细网格 RMSE<1e-2 & Re=\allowbreak{}1000 n=\allowbreak{}33：res=\allowbreak{}9.89e-08，RMSE=\allowbreak{}0.2413；Re=\allowbreak{}100 n=\allowbreak{}49：res=\allowbreak{}9.91e-08，RMSE=\allowbreak{}0.0060 \checkmark \\
A & Re=100 的 Ghia RMSE 网格收敛 & RMSE 单调下降；log(\allowbreak{}RMSE) vs log(\allowbreak{}h) 的斜率在 (\allowbreak{}0.5,\allowbreak{}2.5)(\allowbreak{}一阶迎风约 1) & n=\allowbreak{}17: 0.02024；n=\allowbreak{}25: 0.01211；n=\allowbreak{}33: 0.00870；n=\allowbreak{}49: 0.00603；实测阶 1.10 \checkmark \\
B & Re=50/400 的 dt 限制归属与交叉点 & Re\_cross=\allowbreak{}102.4；Re=\allowbreak{}50 属扩散侧、\allowbreak{}Re=\allowbreak{}400 属 CFL 侧 & Re=\allowbreak{}50: dt\_diff=\allowbreak{}0.01221,\allowbreak{} dt\_cfl=\allowbreak{}0.02500,\allowbreak{} binding=\allowbreak{}0；Re=\allowbreak{}400: dt\_diff=\allowbreak{}0.09766,\allowbreak{} dt\_cfl=\allowbreak{}0.02500,\allowbreak{} binding=\allowbreak{}1 \checkmark \\
B & 中心格式阈值与格点雷诺数 & Re\_th=\allowbreak{}2(\allowbreak{}n-1)；Re=\allowbreak{}1000 在 n=\allowbreak{}33/\allowbreak{}41/\allowbreak{}49 上都已超过阈值；阈值处 Re\_h$\approx$2 & n=\allowbreak{}33: Re\_th=\allowbreak{}64,\allowbreak{} Re\_h(\allowbreak{}Re=\allowbreak{}1000)=\allowbreak{}31.25,\allowbreak{} Re\_h(\allowbreak{}Re\_th)=\allowbreak{}2.000；n=\allowbreak{}41: Re\_th=\allowbreak{}80,\allowbreak{} Re\_h(\allowbreak{}Re=\allowbreak{}1000)=\allowbreak{}25.00,\allowbreak{} Re\_h(\allowbreak{}Re\_th)=\allowbreak{}2.000；n=\allowbreak{}49: Re\_th=\allowbreak{}96,\allowbreak{} Re\_h(\allowbreak{}Re=\allowbreak{}1000)=\allowbreak{}20.83,\allowbreak{} Re\_h(\allowbreak{}Re\_th)=\allowbreak{}2.000 \checkmark \\
A & 四页端点帧的有限性、JSON 常量与 requiredAux & 非有限常量 =\allowbreak{} 0；帧数 =\allowbreak{} 6；每页 requiredAux 都在数组名里 & 6 帧、\allowbreak{}58 个数组，非有限常量计数 0(\allowbreak{}cavity\_vorticity: 10 数组 /\allowbreak{} 33 meta；ghia\_benchmark: 14 数组 /\allowbreak{} 18 meta；turing\_pattern: 7 数组 /\allowbreak{} 34 meta；dam\_break: 10 数组 /\allowbreak{} 20 meta) \checkmark \\
A & 方腔高/低 Re 的残差、RMSE 与时间步绑定 & 高 Re：残差 <=\allowbreak{}1e-6、\allowbreak{}RMSE>0.15、\allowbreak{}CFL 绑定、\allowbreak{}Re\_h>2；低 Re：RMSE<0.02；residual[-1]=\allowbreak{}log10(\allowbreak{}final) & Re=\allowbreak{}1000/\allowbreak{}n=\allowbreak{}41：残差 1.00e-07，RMSE 0.2271，binding=\allowbreak{}1，Re\_h=\allowbreak{}25.0；Re=\allowbreak{}100/\allowbreak{}n=\allowbreak{}49：RMSE 0.0060；residual[-1]=\allowbreak{}-7.000014 vs log10(\allowbreak{}final)=\allowbreak{}-7.000014；中心阈值=\allowbreak{}80.0 \checkmark \\
A & Ghia 四网格 RMSE 曲线与抄写地板 & n\_grid=\allowbreak{}[17,\allowbreak{}25,\allowbreak{}33,\allowbreak{}49]；RMSE 严格递减；斜率$\in$[1.0,\allowbreak{}1.2]；n=\allowbreak{}25 的 RMSE >=\allowbreak{}1.8$\times$n=\allowbreak{}49；中心距 <0.05 & n=\allowbreak{}[17,\allowbreak{} 25,\allowbreak{} 33,\allowbreak{} 49]，RMSE=\allowbreak{}[0.0202 0.0121 0.0087 0.006 ]；斜率=\allowbreak{}1.105，best\_n=\allowbreak{}49；n=\allowbreak{}25 RMSE=\allowbreak{}0.0121 vs n=\allowbreak{}49 0.0060；抄写地板=\allowbreak{}1.0e-02，中心距=\allowbreak{}0.0174；基准点 17 个，u 数组同长=\allowbreak{}True \checkmark \\
A & 四预设的 turing\_\allowbreak{}band / base\_\allowbreak{}valid / $\lambda$(0) / 波长有效性 & turing：band=\allowbreak{}1、\allowbreak{}base=\allowbreak{}1、\allowbreak{}dt\_clamped=\allowbreak{}1、\allowbreak{}波长差<12\%；stripes：band=\allowbreak{}0、\allowbreak{}base=\allowbreak{}1、\allowbreak{}$\lambda$(\allowbreak{}0)>0；spots：band=\allowbreak{}0、\allowbreak{}base=\allowbreak{}0；uniform：band=\allowbreak{}0、\allowbreak{}wavelength\_valid=\allowbreak{}0 且 wavelength=\allowbreak{}0 & turing：$\varphi$*=\allowbreak{}53.899，$\lambda$*=\allowbreak{}0.01502，band=\allowbreak{}1，dt 0.25->0.09537(\allowbreak{}limit 0.09537)，波长 0.1253 vs 0.1166(\allowbreak{}差 7.5\%)；stripes：band=\allowbreak{}0，base=\allowbreak{}1，$\lambda$(\allowbreak{}0)=\allowbreak{}+0.00539；spots：band=\allowbreak{}0，base=\allowbreak{}0；uniform：band=\allowbreak{}0，wavelength\_valid=\allowbreak{}0，wavelength=\allowbreak{}0.000 \checkmark \\
A & 干底溃坝的零步帧、扇区 L1、质量漂移与 CFL & t=\allowbreak{}0：n\_steps=\allowbreak{}0、\allowbreak{}L1=\allowbreak{}0、\allowbreak{}h=\allowbreak{}=\allowbreak{}h\_initial；t=\allowbreak{}0.06：l1\_all>l1\_fan>0、\allowbreak{}0<mass\_drift<1e-6、\allowbreak{}max\_cfl<=\allowbreak{}0.8；n=\allowbreak{}100$\to$400 扇内 L1 递减 & t=\allowbreak{}0：n\_steps=\allowbreak{}0，L1=\allowbreak{}0.0e+00，h 逐位相等=\allowbreak{}True；t=\allowbreak{}0.06：fan=\allowbreak{}113 点，valid=\allowbreak{}166 点，L1\_fan=\allowbreak{}2.0200e-02，L1\_all=\allowbreak{}2.2075e-02，比值=\allowbreak{}1.093；质量漂移=\allowbreak{}9.90e-09，max\_cfl=\allowbreak{}0.800；n=\allowbreak{}100/\allowbreak{}400 扇内 L1=\allowbreak{}3.0410e-02/\allowbreak{}1.2489e-02 \checkmark \\
B & 四页数组与直接调用 numcore 的最大差 & 方腔场/\allowbreak{}Ghia 基准/\allowbreak{}图灵 v/\allowbreak{}溃坝 h 的最大差都 =\allowbreak{} 0 & 方腔 $\omega$ 0.0e+00、\allowbreak{}$\psi$ 0.0e+00、\allowbreak{}x 0.0e+00；Ghia u\_num 0.0e+00、\allowbreak{}u\_ref 0.0e+00；图灵 v 0.0e+00；溃坝 h 0.0e+00、\allowbreak{}x 0.0e+00 \checkmark \\
A & Ghia 三列的表长与端点 & 每列 17 点；u(\allowbreak{}0)=\allowbreak{}0、\allowbreak{}u(\allowbreak{}1)=\allowbreak{}1、\allowbreak{}v(\allowbreak{}0)=\allowbreak{}v(\allowbreak{}1)=\allowbreak{}0；采样严格递增 & Re=\allowbreak{}100: u$\in$[-0.21090,\allowbreak{}1.00000]，v$\in$[-0.24533,\allowbreak{}0.17527]；Re=\allowbreak{}400: u$\in$[-0.32726,\allowbreak{}1.00000]，v$\in$[-0.44993,\allowbreak{}0.30203]；Re=\allowbreak{}1000: u$\in$[-0.38289,\allowbreak{}1.00000]，v$\in$[-0.51550,\allowbreak{}0.37095] \checkmark \\
A & 主涡位置随 Re 的变化与转载分歧地板 & x、\allowbreak{}y 随 Re 单调减小；Re=\allowbreak{}1000 存在两个底角二次涡；分歧地板 1e-2 & Re=\allowbreak{}100 主涡 (\allowbreak{}0.6172,\allowbreak{} 0.7344)，Re=\allowbreak{}400 (\allowbreak{}0.5547,\allowbreak{} 0.6055)，Re=\allowbreak{}1000 (\allowbreak{}0.5313,\allowbreak{} 0.5625)；二次涡 ['bottom\_left',\allowbreak{} 'bottom\_right']；floor=\allowbreak{}0.01 \checkmark \\
A & 图灵底态的两支与稳定性 & u\_-=\allowbreak{}(\allowbreak{}0.25,\allowbreak{}0.3)，det J=\allowbreak{}+0.0045；u\_+ 分支 det<0；spots 无非平凡稳态 & u\_-=\allowbreak{}(\allowbreak{}0.250000,\allowbreak{}0.300000),\allowbreak{} det=\allowbreak{}0.004500；u\_+=\allowbreak{}(\allowbreak{}0.750000,\allowbreak{}0.100000),\allowbreak{} det=\allowbreak{}-0.001500；spots base=\allowbreak{}None \checkmark \\
A & 最不稳定波数与 det 顶点 & $\varphi$*=\allowbreak{}53.899、\allowbreak{}$\lambda$*=\allowbreak{}0.015017；det 顶点 $\varphi$=\allowbreak{}62.75、\allowbreak{}$\lambda$=\allowbreak{}0.013392；经典三预设最快模都在 $\varphi$=\allowbreak{}0，其中 stripes 的 $\lambda$(\allowbreak{}0)=\allowbreak{}+0.0054 是均匀失稳，不是图灵带 & $\varphi$*=\allowbreak{}53.898800，$\lambda$*=\allowbreak{}0.0150170；det 顶点 $\varphi$=\allowbreak{}62.7495，$\lambda$=\allowbreak{}0.0133924，低 10.8\%；经典预设 (\allowbreak{}0.00,\allowbreak{}-0.0350)、\allowbreak{}(\allowbreak{}0.00,\allowbreak{}0.0054)、\allowbreak{}(\allowbreak{}0.00,\allowbreak{}-0.0140)；stripes $\lambda$(\allowbreak{}0)=\allowbreak{}+0.0054 \checkmark \\
B & dt=0.25 在 n=128 的 turing 预设上被夹紧 & limit=\allowbreak{}0.0953674；实际 dt=\allowbreak{}limit，dt\_requested=\allowbreak{}0.25；clamp=\allowbreak{}False 抛 ValueError & limit=\allowbreak{}0.0953674，dt=\allowbreak{}0.0953674，clamped=\allowbreak{}True，requested=\allowbreak{}0.25 \checkmark \\
B & 五步内 $\Sigma$u 变化与 dt$\cdot$$\Sigma$反应 的逐位对账 & 每步差 <=\allowbreak{} 1e-12(\allowbreak{}相对)；扩散项的周期和恒为 0 & predicted=\allowbreak{}761.520552420045，actual=\allowbreak{}761.520552420045，相对差=\allowbreak{}0.00e+00；逐步 ['1.14e-13',\allowbreak{} '1.14e-13',\allowbreak{} '1.14e-13',\allowbreak{} '0.00e+00',\allowbreak{} '0.00e+00'] \checkmark \\
A & 五点模板增长率误差的网格收敛 & 误差按 n\textasciicircum{}-2 下降；n=\allowbreak{}64/\allowbreak{}128/\allowbreak{}256/\allowbreak{}512 的 log-log 斜率约 -2 & 64:9.32e-03、\allowbreak{}128:3.24e-03、\allowbreak{}256:8.68e-04、\allowbreak{}512:2.21e-04；斜率 -1.81 \checkmark \\
B & 经典 spots 的 $\lambda$* 与饱和振幅 & $\lambda$*=\allowbreak{}-0.035<0，但 4000 步后振幅>0.5：自复制瞬态，不是图灵失稳 & $\varphi$*=\allowbreak{}0.000，$\lambda$*=\allowbreak{}-0.0350，振幅=\allowbreak{}0.729，斑点=\allowbreak{}16，t=\allowbreak{}1000.0 \checkmark \\
A & 图灵饱和图案的波长与振幅 & 理论波长 2$\pi$/\allowbreak{}$\varphi$*=\allowbreak{}0.1166；实测应在 12\% 内，振幅>0.2，局部极大>=\allowbreak{}40 & 实测波长=\allowbreak{}0.1253，理论=\allowbreak{}0.1166，相对差=\allowbreak{}7.5\%；振幅=\allowbreak{}0.302，斑点=\allowbreak{}96 \checkmark \\
A & Rusanov 一致性、波速与 t=0 极限 & 相同态通量逐位相等；a=\allowbreak{}|u|+sqrt(\allowbreak{}gh)；t=\allowbreak{}0 时 h=\allowbreak{}1[x<0.5]，u=\allowbreak{}0 & max|$\Delta$F\_h|=\allowbreak{}0.0e+00，max|$\Delta$F\_hu|=\allowbreak{}0.0e+00，波速 [2.9264282071829006,\allowbreak{} 3.681034829318991,\allowbreak{} 4.6794469180700204] \checkmark \\
A & Ritter 干底解的扇区公式与正性 & 扇区公式逐点成立；扇外 h=\allowbreak{}1/\allowbreak{}0；h>=\allowbreak{}0 & 扇内 max|$\Delta$h|=\allowbreak{}0.00e+00，max|$\Delta$u|=\allowbreak{}0.00e+00；h\_min=\allowbreak{}0.0e+00 \checkmark \\
A & 干底溃坝的窗口误差与阶 & 误差单调下降；扇区阶约 0.7-0.9；全域/\allowbreak{}扇区比 >1.3(\allowbreak{}干锋误差) & n=\allowbreak{}100: tight=\allowbreak{}1.91e-02,\allowbreak{} all=\allowbreak{}3.13e-02,\allowbreak{} CFL=\allowbreak{}0.80；n=\allowbreak{}200: tight=\allowbreak{}1.20e-02,\allowbreak{} all=\allowbreak{}2.07e-02,\allowbreak{} CFL=\allowbreak{}0.80；n=\allowbreak{}400: tight=\allowbreak{}7.42e-03,\allowbreak{} all=\allowbreak{}1.30e-02,\allowbreak{} CFL=\allowbreak{}0.80；阶 [0.6691266073450176,\allowbreak{} 0.6942551943035056]；全域/\allowbreak{}扇区=\allowbreak{}1.75 \checkmark \\
B & 周期湿域与干底的质量漂移 & 周期：<=\allowbreak{}1e-12；干底：>1e-10(\allowbreak{}正性地板注入，不是不守恒) & periodic drift=\allowbreak{}2.22e-16；dry drift=\allowbreak{}9.90e-09 \checkmark \\
B & 静水漂移：well-balanced vs 朴素源项 & well-balanced <1e-14；naive >1e-4；两者应差多个数量级 & well-balanced=\allowbreak{}1.89e-16，naive=\allowbreak{}3.30e-04，比值 1.7e+12 \checkmark \\
B & CFL 稳定/失稳的实测边界 & cfl=\allowbreak{}0.8：ht\_end 到达、\allowbreak{}h\_min>0.2；cfl=\allowbreak{}2.0：非有限或 h\_max<1e-6 & stable: reached=\allowbreak{}True,\allowbreak{} h\_min=\allowbreak{}0.426,\allowbreak{} maxCFL=\allowbreak{}0.80；bad: finite=\allowbreak{}False,\allowbreak{} h\_max=\allowbreak{}inf,\allowbreak{} steps=\allowbreak{}830 \checkmark \\
A & 二维径向溃坝的有限性、正性与出流质量 & h\_min>0；出流边界质量/\allowbreak{}能量不增(\allowbreak{}允许舍入 1e-6)；推到 t\_end & t=\allowbreak{}0.0500,\allowbreak{} steps=\allowbreak{}40,\allowbreak{} h\_min=\allowbreak{}1.000e-01,\allowbreak{} mass=\allowbreak{}0.278418 (\allowbreak{}m0=\allowbreak{}0.278418),\allowbreak{} energy=\allowbreak{}0.852330 (\allowbreak{}e0=\allowbreak{}1.011704),\allowbreak{} maxCFL=\allowbreak{}0.30 \checkmark \\
\end{longtable}

